\documentclass[11pt,a4paper]{article}
\usepackage[T1]{fontenc}
\usepackage[utf8]{inputenc}
\usepackage[ngerman]{babel}
\usepackage[margin=2.5cm]{geometry}
\usepackage[hidelinks]{hyperref}
\usepackage{array}
\setlength{\parindent}{0pt}
\setlength{\emergencystretch}{2em}
\renewcommand{\arraystretch}{1.25}
\sloppy
\begin{document}
\shorthandoff{"}
\begin{center}
  {\Large\textbf{Webtechnologien -- Übungsblatt 5}}
\end{center}
\section*{Aufgabe 4: Retrieving Products via XMLHttpRequest}
\textbf{Aufgabe:} In \texttt{products.js} die Methode
\texttt{initialize(jsonURL, domElement)} schreiben. Sie lädt
\texttt{products.json} mit \texttt{XMLHttpRequest} vom Webserver und
zeigt jedes Produkt im übergebenen \texttt{domElement} so an wie auf dem
letzten Blatt.

\medskip
\textbf{Ergebnis:} \texttt{products.js} ist die Originaldatei. Was wegfällt,
ist nur auskommentiert; alles Neue steht nach fünf Leerzeilen am
Dateiende zwischen \texttt{[NEU]} und \texttt{[ENDE NEU]}.
\begin{itemize}
  \item \texttt{[4 GEÄNDERT]}: Die URL \texttt{http://localhost:8080/...}
    passt nicht, weil Apache von XAMPP auf Port 80 läuft. Neu ist
    \texttt{let productURL = "products.json"}: eine relative URL, also
    \glqq gleicher Ordner wie \texttt{products.html}\grqq. Die passt für
    jeden Ordner und jeden Port.
  \item \texttt{[4 VERSCHOBEN]}: Der Aufruf
    \texttt{initialize(productURL, container)} steht jetzt unten, weil
    \texttt{productURL} erst dort angelegt wird; eine
    \texttt{let}-Variable darf man erst nach ihrer Zeile benutzen.
  \item \texttt{[4 ENTFERNT]}: die leere Vorlage von \texttt{initialize}.
  \item \texttt{initialize}: \texttt{new XMLHttpRequest()} legt die
    Anfrage an, \texttt{xhr.open("GET", jsonURL)} sagt, was geholt wird,
    \texttt{xhr.send()} schickt sie ab. Der Code läuft danach sofort
    weiter; die Antwort kommt später (asynchron) im Event
    \texttt{onload}. Dort wird geprüft, ob \texttt{xhr.status} 200 (OK)
    ist, dann macht \texttt{JSON.parse(xhr.responseText)} aus dem Text
    ein Array mit vier Produkt-Objekten.
  \item \texttt{renderProducts}: baut pro Produkt mit
    \texttt{createElement} dasselbe HTML wie in \texttt{productsA.html}
    von Blatt 4: \texttt{div.product} mit \texttt{data-prodId},
    \texttt{h2} mit Titel, darunter \texttt{div.desc},
    \texttt{div.price} (2890 wird zu \glqq 2890,00\grqq) und
    \texttt{div.buy} mit dem Button. Text wird mit
    \texttt{textContent} gesetzt, also nie als HTML ausgewertet.
  \item Fehler: Bei einem anderen Status (z.\,B. 404) oder kaputtem JSON
    (\texttt{try}/\texttt{catch} um \texttt{JSON.parse}) und bei
    \texttt{onerror} (gar keine Antwort, z.\,B. Server aus) zeigt
    \texttt{showError} eine rote Meldung statt der Produkte.
  \item \texttt{products.html}: Nur ein \texttt{<style>}-Block ist neu
    (\texttt{[NEU]}), kopiert aus \texttt{productsA.html}, damit die
    Produkte gleich aussehen wie auf Blatt 4.
\end{itemize}

\textbf{Warum ein Webserver?} Wird die Seite per Doppelklick über
\texttt{file://} geöffnet, blockiert der Browser aus Sicherheitsgründen
die Anfrage (Konsole: \glqq Access to XMLHttpRequest at
\texttt{file:///...} has been blocked by CORS policy\grqq); es kommt
\texttt{onerror} und die Fehlermeldung. Über
\texttt{http://localhost/...} liefert Apache die Datei aus.

\medskip
\textbf{Hinweis:} Der Ordner liegt in XAMPP unter
\texttt{C:\textbackslash xampp\textbackslash htdocs\textbackslash ex4},
aufrufen mit \url{http://localhost/ex4/products.html} (vorher im XAMPP
Control Panel Apache starten). Der Button \glqq add to cart\grqq{} ist
nur für das gleiche Aussehen da; ein Warenkorb ist nicht verlangt.
\texttt{products.json} ist unverändert.

\medskip
\textbf{Test:} In Chrome:
\begin{center}
\begin{tabular}{|>{\raggedright\arraybackslash}p{5.2cm}|>{\raggedright\arraybackslash}p{8.5cm}|}
\hline
\textbf{Fall} & \textbf{Ergebnis} \\ \hline
\url{http://localhost/ex4/products.html} (XAMPP) & 4 Produkte wie auf
  Blatt 4: Sony Alpha IV 2890,00, Sony 24-105 1300,00, Nicon 900,00,
  Tamron 800,00 \\ \hline
Datei nicht gefunden (404) & \glqq Products could not be loaded: HTTP 404
  Not Found\grqq \\ \hline
kaputtes JSON & Meldung \glqq invalid JSON (...)\grqq \\ \hline
Server nicht erreichbar & Meldung \glqq network error\grqq \\ \hline
über \texttt{file://} geöffnet & Meldung \glqq network error\grqq{} \\ \hline
\end{tabular}
\end{center}
Kein unbehandelter Fehler: In der Konsole stehen nur die Meldung des
Browsers (z.\,B. \glqq Failed to load resource\grqq{} bzw. bei
\texttt{file://} die CORS-Sperre) und die eigene Meldung von
\texttt{showError}.

\medskip
{\small \textbf{Hilfsmittel:} Claude Code (Claude Opus 5.5) für Erstellung, Prüfung und Dokumentation.}
\end{document}
